\chapter{Scientific problem}
\label{chapter:problem}

\section{Informal description}
\label{section:informal}

An auditor receives two artefacts: a policy text and a software repository. The policy is long, hierarchical (articles,
paragraphs, points) and mixes obligations with definitions and procedural provisions. The repository contains source code,
infrastructure-as-code (Terraform, Kubernetes), configuration, documentation and possibly model cards and sample logs. The
auditor must (1) find the obligations that are verifiable in a repository, (2) decide which of them apply to this system
(this depends on the AI Act risk tier, which itself must be inferred from the intended purpose), (3) for each applicable
obligation, locate evidence and decide whether it is met fully, partially or not at all, and (4) explain the decision with
citations so that a third party can check it. The answer ``I cannot tell, a human must look'' is legitimate and preferable
to a wrong answer.

\paragraph{Why an intelligent method.} The mapping from legal text to code is not syntactic: Art.~12 ``automatic recording
of events'' may be implemented as \texttt{logger.info(...)}, as a custom \texttt{DecisionJournal.record(...)} writing to
object storage, or not at all despite a README claiming it. Rules alone suffer from false negatives (indirect evidence)
and false positives (decoys such as commented code, TODOs, marketing claims). Reading and weighing heterogeneous evidence
against an open-textured statement is a natural-language inference task, for which LLMs are the current best tool; their
weaknesses (hallucination, sensitivity to context, cost) are what the rest of the design compensates for.

\paragraph{Advantages and disadvantages of the candidate methods.}
\begin{itemize}[leftmargin=*]
  \item \emph{Rule-based (policy-as-code)}: deterministic, explainable, cheap, CI-friendly; brittle, high maintenance,
        cannot judge sufficiency or proportionality.
  \item \emph{LLM-only}: flexible, reads documents and code alike, writes good rationales; non-deterministic, may
        hallucinate evidence, costs money and time per call, sensitive to retrieval quality~\cite{liu2024lost}.
  \item \emph{Hybrid with verification and abstention} (ours): spends LLM calls only where rules are weak, constrains the
        LLM to cited evidence and abstains when unsure; more moving parts, routing thresholds must be tuned.
\end{itemize}

\begin{figure}[htbp]
\centering
\resizebox{\linewidth}{!}{\graphicalabstract}
\caption{Graphical abstract: policies become requirements, repositories become evidence; a hybrid rule/LLM evaluation with
a citation verifier and a human-in-the-loop gate produces an auditable report.}
\label{fig:graphical-abstract}
\end{figure}

\section{Formal description}
\label{section:formal}

\paragraph{Policy side.} A policy $P$ is a finite tree of clauses $\Cl(P)=\{c_1,\dots,c_m\}$; each clause has a stable
identifier (e.g.\ \texttt{eu-ai-act/Art.12(1)}), a text and a parent. A \emph{decomposition} $\delta:\Cl(P)\to 2^{\Req}$
produces requirements
\[
  r=(\mathrm{id},\ C_r\subseteq\Cl(P),\ \kappa_r,\ s_r,\ A_r\subseteq\Lv,\ w_r,\ K_r,\ Q_r,\ g_r),
\]
where $\kappa_r$ is a category (record keeping, transparency, human oversight, \dots), $s_r$ a normative testable
statement, $A_r$ the risk tiers it applies to with $\Lv=\{\text{prohibited},\text{high},\text{limited},\text{minimal}\}$,
$w_r\in\{\text{critical},\text{high},\text{medium},\text{low}\}$ a severity, $K_r$ a (possibly empty) set of deterministic
checks, $Q_r$ retrieval queries and $g_r$ guidance for the judge. $\delta$ is either curated by hand or computed by an
LLM (PolicyAnalyst).

\paragraph{Repository side.} A repository $D$ is a set of text files. A collector $\varphi$ maps it to an evidence set
$\Ev=\varphi(D)$; each $e\in\Ev$ is a tuple $(\mathrm{id}, k_e, \mathrm{path}, [\ell_1,\ell_2], x_e, \sigma_e)$ with a kind
$k_e\in\{\text{code},\text{config},\text{iac},\text{doc},\text{model\_card},\text{log}\}$, a line span, a snippet $x_e$ and
a set of regex signals $\sigma_e$ (logging call, AI disclosure, human override, retention, \dots). A risk classifier
$\rho:2^{\Ev}\to\Lv\cup\{\text{unknown}\}$ assigns the AI Act tier.

\paragraph{Audit function.} Let $\St=\{\stat{c},\stat{p},\stat{nc},\stat{nr},\stat{na}\}$ (compliant, partial,
non-compliant, needs review, not applicable). The auditor is a function
\[
  f:\ \Req\times 2^{\Ev}\ \longrightarrow\ \St\times[0,1]\times 2^{\Ev},\qquad
  f(r,\Ev)=(\hat y_r,\ \pi_r,\ \hat E_r),
\]
returning a status, a confidence and a set of cited evidence, subject to three hard constraints:
\begin{align}
  &\text{(applicability)} && \rho(\Ev)\notin A_r \;\Rightarrow\; \hat y_r=\stat{na}, \label{eq:applic}\\
  &\text{(grounding)}     && \hat E_r\subseteq \Ev_r \subseteq \Ev, \text{ where } \Ev_r \text{ is the evidence shown to the evaluator}, \label{eq:ground}\\
  &\text{(no unsupported compliance)} && \hat y_r\in\{\stat{c},\stat{p}\}\;\Rightarrow\;\hat E_r\neq\emptyset
     \text{ and } \hat E_r \text{ supports } s_r. \label{eq:support}
\end{align}
The status \stat{nr} is an \emph{abstention} (selective prediction): it is never a ground-truth label, it transfers the
decision to a human at a cost $\lambda$.

\paragraph{Objective.} With human labels $y_r\in\St\setminus\{\stat{nr}\}$, a misclassification cost matrix
$\Lambda$ (asymmetric: predicting \stat{c} when the truth is \stat{nc} or \stat{p} -- a \emph{false compliance} -- is the
most expensive error, because it hides a legal gap) and a per-call LLM cost, the design goal is
\[
  \min_{f}\ \sum_{r}\Big[\Lambda(y_r,\hat y_r)\,\mathbb{1}[\hat y_r\neq\stat{nr}] + \lambda\,\mathbb{1}[\hat y_r=\stat{nr}]\Big]
  \;+\;\mu\cdot\mathrm{Cost}_{\text{LLM}}(f)
  \quad\text{s.t. (\ref{eq:applic})--(\ref{eq:support}).}
\]
We do not optimise this objective numerically (there is no training set large enough); it motivates the routing and
abstention thresholds of Chapter~\ref{chapter:approach} and the metrics below.

\paragraph{Metrics.} Let $\mathcal{A}=\{r: y_r\neq\stat{na}\}$ be the applicable requirements and
$G=\{r\in\mathcal{A}: y_r\in\{\stat{nc},\stat{p}\}\}$ the true \emph{gaps}.
Let $F=\{\stat{nc},\stat{p},\stat{nr}\}$ be the predictions that \emph{flag} a requirement for attention. The metrics
(implemented in \texttt{app/benchmark/evaluate.py}) are:
\begin{itemize}[leftmargin=*]
  \item \textbf{Accuracy} $=\frac{1}{|\mathcal{A}|}\sum_{r\in\mathcal{A}} \mathbb{1}[\hat y_r=y_r]$ on the applicable
        labels; an \stat{nr} or a wrong \stat{na} prediction counts as an error.
  \item \textbf{Macro-F1} over the graded classes $\{\stat{c},\stat{p},\stat{nc}\}$ on $\mathcal{A}$.
  \item \textbf{Gap detection} (binary): gap precision $=|\{r\in G:\hat y_r\in F\}|/|\{r\in\mathcal{A}:\hat y_r\in F\}|$,
        gap recall $=|\{r\in G:\hat y_r\in F\}|/|G|$ and their F1.
  \item \textbf{False-compliance rate} $\mathrm{FCR}=|\{r\in G:\hat y_r=\stat{c}\}|/|G|$ (lower is better; the
        safety-critical failure of an auditor).
  \item \textbf{Needs-review rate} $=|\{r\in\mathcal{A}:\hat y_r=\stat{nr}\}|/|\mathcal{A}|$ (human workload) and
        \textbf{accuracy@deferral}, the accuracy on the applicable items that were not deferred (selective accuracy).
  \item \textbf{Applicability accuracy} $=\frac{1}{|\Req|}\sum_r \mathbb{1}\big[(\hat y_r=\stat{na})=(y_r=\stat{na})\big]$
        over all labels, and \textbf{risk-tier accuracy} of $\rho$ over the applications.
  \item \textbf{Citation validity}: share of cited evidence ids that exist in the evidence index, and share of
        \stat{c} verdicts carrying at least one citation~\cite{rashkin2023ais,gao2023alce}.
  \item \textbf{Complexity}: wall-clock latency per audit and per stage, number of LLM calls, input/output tokens,
        equivalent USD cost reported by the CLI, cache hit rate.
\end{itemize}

\paragraph{Computational complexity.} With $n=|\Ev|$ chunks, $|\Req|$ requirements and $q$ queries per requirement,
evidence collection is linear in the repository size, BM25 indexing is $O(n\bar{\ell})$ and each query $O(n)$; the rule
engine is linear in the scanned text per check. The dominant cost is the number of LLM calls: $|\Req_{\mathcal{A}}|$
judge calls in LLM mode versus $|\{r:\text{route}(r)=\text{LLM}\}|$ in hybrid mode (plus one risk call, and optionally
one verifier call per positive verdict), each with a bounded prompt of at most 14 evidence items / 12\,000 characters.
